\section{Einleitung}
%%% HINWEISE HINNEBURG %%%
% kurze Einleitung --> Zielproblem beschreiben -> Fragestellung motivieren --> mittels Techniken der Informationsvisualisierung beantworten
% Wichtig: Lösung soll aus Argumentation hergeleitet werden. Schlecht: "Das ist das Problem und ich zeichne diese 3 Diagramme."
\textbf{...TODO...}

\subsection{Anwendungshintergrund}
% Warum ist es nicht trivial, diese Fragestellung zu beantworten?
% Sie müssen genug Hintergrund bereitstellen, so dass die Lesenden sich ein Urteil bilden können, ob ihre Lösung funktioniert. Sie sollen die Lesenden jedoch nicht mit Anwendungsdetails so überschütten, dass der Fokus auf die Fragen zur Informationsvisualisierung untergehen. 
\textbf{...TODO...}

\subsection{Zielgruppen}
%Beschreiben Sie die Personengruppe oder Personengruppen, die das von Ihnen benannte Anwendungsproblem lösen möchte. Auf welches Vorwissen können Sie in dieser Gruppe von Anwenderinnen aufbauen? Welche Informations"-bedürf"-nisse werden durch die Visualisierungen adressiert?
\textbf{...TODO...}

\subsection{Überblick und Beiträge}
%In diesem Abschnitt geben sie einen kurzen Überblick über die Daten und verwendeten Visualisierungen.
%Dann benennen sie die Beiträge ihres Projekts. Diese Beiträge müssen sie in den hinteren Teilen des Berichts genauer ausführen und belegen.
% --> Erst zum Schluss schreiben. z.B. ... Durch unsere Lösung sind die Daten in eiem Diagramm übersichtlich darstellbar ...
% --> vorgreifen auf Zusammenhang zwischen den Visualisierungen
\textbf{...TODO...}
